\begin{abstract}
Первый том не объявляет завершённую физическую теорию. Его задача ---
зафиксировать исходную исследовательскую программу, ввести единый язык
описания и определить, каким требованиям должна удовлетворять будущая
проверяемая модель. Геометрические, спектральные, топологические, фазовые и
корреляционные конструкции рассматриваются как разные способы чтения одного
возможного структурного источника. Том устанавливает границу между рабочей
онтологией, полуформальным каркасом и доказанным физическим результатом.
\end{abstract}

\chapter{Назначение первого тома}

Исследование начинается не с готовой системы аксиом, а с повторяющегося
структурного наблюдения. Разные математические языки проекта снова и снова
выделяют глобальную допустимость, спектральные ограничения, топологические
классы и устойчивые режимы. Такое совпадение ещё не является доказательством
единой теории, но задаёт предмет строгой проверки.

Основная постановка состоит в следующем: частицы, поля, заряды и геометрия
могут быть не независимыми первичными объектами, а устойчивыми режимами более
глубокой структуры допустимых состояний и переходов. Поэтому первый том должен
не подгонять известные числа, а сформулировать условия, при которых эта мысль
становится научной программой.

\chapter{Рабочая онтология}

\section{Первичность структуры}

Первичным кандидатом считается не набор готовых вещей, а закон согласования
состояний, переходов и ограничений. Частица в такой картине является устойчивой
формой организации, поле --- способом связывания состояний, а геометрия ---
частью условий их существования.

Это утверждение имеет статус рабочей гипотезы. Оно не заменяет операторного
построения, уравнений движения и наблюдательной проверки.

\section{Допустимость раньше локальности}

Не всякая локально записанная конфигурация допустима глобально. Она должна
быть совместима с топологией, спектром, спиновой структурой, граничными
условиями и голономией. Поэтому пространство допустимых режимов логически
предшествует их локальному эффективному описанию.

\section{Несколько языков одного объекта}

В проекте различаются три основных языка:
\begin{itemize}
 \item геометрический язык описывает носители, кривизну, объёмы и циклы;
 \item спектральный язык описывает операторы, собственные значения, следы и
       асимптотические коэффициенты;
 \item корреляционный язык описывает согласование локальных и глобальных
       частей как проявлений одной структуры.
\end{itemize}

Ни один язык не объявляется достаточным заранее. Рабочая теория должна
показать, какие данные сохраняются при переходе между ними и какие сведения
не восстанавливаются однозначно.

\chapter{Гипотеза единого источника}

Предполагается существование структуры $\mathfrak S$, допускающей
геометрическое, спектральное и корреляционное чтения. Наблюдаемые физические
режимы должны возникать как её согласованные редукции.

Минимальные требования к $\mathfrak S$ таковы:
\begin{enumerate}
 \item она задаёт пространство допустимых состояний;
 \item она порождает операторы и спектральные инварианты;
 \item она несёт глобальную топологическую и фазовую информацию;
 \item она допускает переход к локальным наблюдаемым;
 \item она содержит инвариантное ядро, сохраняемое при допустимой смене
       описания.
\end{enumerate}

Гипотеза считается содержательной только тогда, когда эти требования можно
проверить независимо друг от друга и сформулировать условия её отказа.

\chapter{Исследовательский протокол}

Метод S2T переводит обнаруженный структурный след в проверяемую теорию через
последовательность обязательных этапов.

\section{Выделение инвариантного ядра}

Сначала фиксируются утверждения, сохраняющиеся в нескольких языках описания.
Совпадение обозначений или чисел само по себе не считается инвариантом: нужен
явный переход между объектами и контроль его неоднозначности.

\section{Минимальная формализация}

Затем вводится наименьший математический объект, способный нести найденное
ядро. Новые сущности и коэффициенты допускаются только при наличии отдельной
роли и независимого способа проверки.

\section{Вывод редукций}

Геометрическое, спектральное, корреляционное и эффективное чтения должны быть
получены из одного объекта. Если связь вводится отдельным правилом для каждого
сектора, единый источник не считается построенным.

\section{Проверка прежних результатов}

Новая конструкция обязана воспроизводить сохранённые результаты проекта без
скрытых нормировок и послепроверочного выбора ветви. Отдельно проводятся
удаление элементов модели, перестановочный контроль и проверка устойчивости
знаков.

\section{Новые различающие следствия}

Программа достигает физического уровня только после получения следствий,
которые не использовались при построении модели. Для каждого такого следствия
заранее задаются допустимая погрешность и критерий провала.

\chapter{Граница утверждений}

В первом томе необходимо различать четыре статуса:
\begin{itemize}
 \item философская мотивация;
 \item рабочая гипотеза;
 \item математически определённая конструкция;
 \item физически проверенное следствие.
\end{itemize}

Переход между статусами требует отдельного доказательства. Выразительная
интерпретация не заменяет определения, численное совпадение не заменяет вывода,
а существование топологического класса не заменяет устойчивой динамики.

\chapter{Следующий сегмент}

Следующая часть должна определить инвариантное ядро, типы структурных режимов
и правила отнесения наблюдаемых к геометрическому, фазовому, внутреннему или
смешанному сектору.